% =====================================================================
%  Chapter 25 — The \idxw{Array} Architecture and the \idxw{Service} Model
% =====================================================================

\chapter{The Array Architecture and the Service Model}
\label{ch:architecture}

The arithmetic of the previous five chapters constrains the architecture
more tightly than a concept paper would suggest, and the constraint is
productive: it removes the mid-infrared channel, it fixes the baseline
range, it sets the wavelength, and it makes the station a service\idx{Service} rather
than an instrument. This chapter assembles the result.

\section{Four elements}

\subsection{Site and plant}

A sunlit ridge or crater rim at the lunar south pole, within a few hundred
metres of a permanently shadowed region used for services and power
conversion, so that heat rejection and power generation are decoupled from
the aperture geometry. NASA's own lunar surface technology catalogue treats
10 kilowatts of power delivered to permanently shadowed regions as an
available capability.

The optical channel requires only that the baffle and sunshade be held
below about 250 kelvin, which is an insulation problem. A mid-infrared
channel is optional, expensive, and on the evidence of Chapter 23 it loses to
Sun-Earth L1, so the proposed architecture does not carry one.

\subsection{The array}

A distributed set of fixed apertures on baselines of 300 metres to
3 kilometres, each fixed to a surveyed anchor point, with at least two
orthogonal orientations because the baseline orientation sweeps with the
Moon's rotation and a single orientation cannot deliver a two-angle
position.

The baseline is set by geology rather than by metrology. The requirement is
the anchor-point separation and its survey to sub-millimetre over kilometres,
which is a lunar geodesy problem rather than an optical-metrology problem.
The separations have to be discrete and selectable, because
\tabref{tab:array-separations} in Chapter 22 shows the matched baseline
changing with the target's distance by a factor of twelve across the
class.

\subsection{Wavelength and channels}

Visible, centred near 750 nanometres, which is where the 50 to 140 metre
class is brightest --- apparent magnitudes from 12.4 to 25.8 across the class
and the 0.1 to 1.3 astronomical unit range --- and where the concept study's
own sensitivity analysis sits.

The fringe visibility in the resolved regime is set by the body's limb
profile and surface contrast rather than by a uniform disc, which is why the
literature finds that a Minnaert model with multiple baselines is necessary
for a few per cent on size. The book's own generator uses the uniform-disc
model and says so in Chapter 20; the requirement it derives is insensitive
to the model, and the achievable precision is not.

A photometric and low-resolution spectroscopic channel on the same apertures
supplies the albedo prior and the taxonomy the thermal solution requires,
and the lightcurve that fixes the rotation period.

\subsection{Service}

This is the element with the highest verified value and the lowest technical
difficulty, and it is where the owner's original concept finds its proper
place.

The station is not a survey. It is a named-target characterisation service\idx{Service}
with a bounded turnaround, and each request returns a resolved shape model,
a pole and period, the epoch-corrected astrometry, and the photometry
needed for the albedo prior.

The requests are generated by an automated alert path\idx{Alert path} from the risk list to
the instrument, because the documented failure of the 2024~YR4 episode was
that ``no automated system currently alerts physical characterization
sites'', and the 2025 experiment with 2025~FA22 showed that such an alert is
operationally simple.

\begin{seamnote}{What the beam to Earth became}
The original concept carried a beam from the Moon to the Earth, and the
instinct behind it was correct: something had to connect the risk list to
an instrument that could respond. The record says the connection was the
missing piece. The beam is therefore re-purposed as an alerting and
scheduling path rather than as a ranging channel, because the ranging
channel turns out to be the weaker of the two claims, and because the
documented failure was a routing failure rather than a sensing failure.
\end{seamnote}

\section{The service rate, and what it means}

The service model has a number attached to it, and the number is what makes
the proposal honest.

A fixed station needs a 90-degree sweep of its baseline orientation to
deliver the two orthogonal measurements an interferometric position
requires, and the Moon's rotation takes 6.83 days to deliver that sweep.
Allowing thirty per cent of the time for overhead --- for the night, for
dust events, for maintenance, for scheduling conflicts --- one station runs
about 37 campaigns a year.

\begin{table}[!htbp]
\centering
\small
\caption{The service rate against the demand, and the coverage it buys.}
\label{tab:service}
\begin{tabularx}{\textwidth}{@{}P{1.56in}P{1.00in}Y@{}}
\toprule
\rowcolor{panel}
\textbf{Quantity} & \textbf{Value} & \textbf{Source} \\
\midrule
Orientation sweep & 6.83 days & the Moon's sidereal rotation \\
\rowcolor{panel}
Campaigns per year & 37 & 30 per cent overhead allowance \\
New 140 m objects per year & 342 & CNEOS, 2025 \\
\rowcolor{panel}
Annual coverage of the arrival rate & 10.9 per cent & the two above \\
Known 140 m objects & 11,725 & CNEOS, 22 November 2026 \\
\rowcolor{panel}
Years to clear the known catalogue & 313 & the two above \\
\bottomrule
\end{tabularx}
\end{table}

\tabref{tab:service} is the number that decides what the station is. It is
not a census machine. It cannot be, and no design change to the optics
changes the arithmetic, because the binding constraint is the rotation rate
and not the instrument.

\begin{deeper}{The rate against the backlog is the wrong comparison, and the
book says so}
Three hundred and thirteen years to clear the known catalogue sounds like a
disqualification, and in a planning document it would be. But the comparison
is wrong, for the reason Chapter 4 established: the gap is not discovery and
the station is not a census. The correct comparison is against the objects
whose characterisation changes an outcome --- the ones that will return, the
ones whose keyhole sits in the tail, the ones whose energy decides a
hazard class. That population is orders of magnitude smaller than 11,725,
and the selection rule of Chapter 30 defines it. A capability that serves
the decisive objects at ten per cent of the annual flow, with a turnaround
of days rather than years, is a triage capability, and triage is what the
gap requires.
\end{deeper}

\section{Why the service model is the load-bearing element}

It is worth being explicit about which of the four elements carries the
argument, because it is not the array.

The array is the technically interesting element and it is the one the
concept literature has already proposed. The wavelength and channels follow
from the photometry. The site and plant follow from the thermal budget.

The service model is the one that makes the whole thing an instrument rather
than an experiment, and it is the one with the smallest technical difficulty
and the largest verified value. The documented failure mode of the 2024~YR4
episode was that the observers of the object did not know it had reached
Torino Level 1. The documented capability gap is that no automated system
alerts physical characterisation sites. The documented demonstration that
the gap is closable is the 2025~FA22 alert exercise.

A station that returns a resolved shape, a pole and a period within days
of a request, triggered automatically by the risk list, is a service. A
station that observes what it can see is an experiment.

\begin{keytakeaways}
\item The architecture is a site, an array of surveyed anchors, a visible
channel with a photometric and spectroscopic companion, and a service.
\item The service returns a resolved shape, a pole and period, the astrometry
and the photometry, on a bounded turnaround, triggered by the risk list.
\item The rate is 37 campaigns a year, 10.9 per cent of the annual arrival
rate of the hazard class, and 313 years to clear the known catalogue.
\item The service model, not the optics, is the element that closes the
documented gap and that the 2025 alert exercise shows to be simple.
\end{keytakeaways}
